\section{Control as inference} \label{sec:control}

Sections~\ref{sec:pf}--\ref{sec:bridge} treat the dynamics parameters
$\thetadyn$ as the unknown and the schedule $U$ as given. This section
reverses the perspective: $\thetadyn$ is held fixed at its posterior
mean (or sampled from the posterior cloud), and the schedule itself
becomes the unknown. The framework solves the resulting stochastic-
optimal-control problem by re-using the SMC$^2$ engine of
Section~\ref{sec:smc2}, with the data-likelihood replaced by an
exponential of the negative cost.

\subsection{The cost functional}

For a fixed dynamics-parameter vector $\thetadyn$, an initial state
$x_0$, and a candidate schedule $U: [0, T] \to \R^{n_u}$, define the
trajectory cost
\begin{equation}
J(U) \;=\; \E_{X \mid U}\!\left[\,
   \int_0^T \ell\bigl(X(t),\, U(t)\bigr)\, \dd t
   \,+\, V_f\bigl(X(T)\bigr)\,\right],
\label{eq:cost}
\end{equation}
where the expectation is over the law of $X$ under the SDE
\eqref{eq:sde} driven by $U$, $\ell$ is a stage cost, and $V_f$ is a
terminal cost. The framework treats $\ell$ and $V_f$ as user-supplied
callables; common choices include negative running reward, soft state
barriers, and effort-quadratic penalties. Concrete examples are given
in Section~\ref{sec:control-cost-variants}.

\subsection{Schedule decoder}

The schedule $U$ is parameterised in finite dimensions by a vector
$\thetactrl \in \R^{d_c}$. The framework's reference decoder is a
Gaussian-RBF basis: with $\{\varphi_a(t)\}_{a=1}^{d_c}$ a fixed basis
on $[0, T]$, a logit bias $c_U$ chosen so that $\thetactrl = 0$
produces a canonical default $U_{\rm def}$, and an upper bound
$U_{\max}$,
\begin{equation}
U(t;\,\thetactrl) \;=\;
   U_{\max} \cdot \sigmoid\!\Biggl(\,c_U
      + \sum_{a=1}^{d_c} (\thetactrl)_a\,\varphi_a(t)\Biggr).
\label{eq:schedule-decoder}
\end{equation}
The sigmoid envelope keeps $U(t; \thetactrl) \in [0, U_{\max}]$ for
every $\thetactrl$; the search is unconstrained in $\R^{d_c}$. Other
decoders (piecewise constant, neural network, etc.) plug in via the
same \texttt{ControlSpec} interface; the framework only requires the
decoder to be a JAX-traceable mapping
$\R^{d_c} \to \R^{n_u \times T}$.

\subsection{Tempered control posterior}

Following the control-as-inference tradition (Kappen, 2005;
Toussaint, 2009; Levine, 2018), the framework defines a probability
measure on $\thetactrl$ by treating the negative cost as an
unnormalised log-likelihood. With inverse temperature $\beta > 0$ and
a Gaussian prior
$\thetactrl \sim \Ncal(0,\, \sigma_{\rm prior}^2 I_{d_c})$, the
\emph{control posterior} is
\begin{equation}
q_\beta(\thetactrl) \;\propto\;
   \exp\bigl(-\beta\,J(\thetactrl)\bigr)\;
   \Ncal(\thetactrl;\,0,\,\sigma_{\rm prior}^2 I_{d_c}),
\label{eq:control-posterior}
\end{equation}
where $J(\thetactrl) := J(U(\,\cdot\,; \thetactrl))$. Sampling from
$q_\beta$ is performed by the SMC$^2$ machinery of
Section~\ref{sec:smc2}, but with the unbiased filter likelihood
$\Lhat_N(\thetadyn)$ replaced by the cost-likelihood
$\exp(-\beta\,J(\thetactrl))$. The inverse temperature $\beta$ is
calibrated automatically from the cost spread of the prior cloud, and
the tempered schedule $\lambda \in [0, 1]$ smoothly interpolates from
prior to control posterior.

\begin{proposition}[Posterior mean as a regularised cost minimiser]
\label{prop:control-mean}
Let $q_\beta$ be the control posterior \eqref{eq:control-posterior}
and write $\bar\theta := \E_{q_\beta}[\thetactrl]$. Then
\begin{equation}
\bar\theta \;=\;
   \argmin_{\theta \in \R^{d_c}}
   \Bigl\{\E_{q_\beta}\bigl[\|\theta - \thetactrl\|_2^2\bigr]\Bigr\},
\label{eq:posterior-mean-l2}
\end{equation}
i.e.\ the posterior mean is the $L^2$ minimiser of squared distance
under the tempered measure. When $J$ is approximately quadratic near
its minimum and the prior is not too informative, $\bar\theta$ sits
in a neighbourhood of the unregularised minimiser
$\argmin_\theta J(\theta)$ of size $O(\beta^{-1/2})$.
\end{proposition}

The framework uses $\bar\theta$, computed as the equally-weighted mean
of the SMC$^2$ parameter cloud at $\lambda = 1$, as the controller
output. The corresponding schedule is then
$U(\,\cdot\,; \bar\theta)$ via \eqref{eq:schedule-decoder}.

\subsection{Hard-indicator vs soft-sigmoid chance-constrained variants}
\label{sec:control-cost-variants}

Many practical control problems are most naturally specified via a
\emph{chance constraint} on a state event rather than via a smooth
running cost. A typical form is
\begin{equation}
\Pcal\bigl[X_t \in \mathcal{B}_{\rm bad}(U_t)\bigr] \;\le\; \alpha,
\qquad \forall\, t \in [0, T],
\label{eq:chance-constraint}
\end{equation}
where $\mathcal{B}_{\rm bad}$ is some user-defined ``bad set'' (e.g.\
the basin of attraction of a collapsed equilibrium). The framework
supports two ways of encoding this constraint inside the
control-posterior cost.

\paragraph{(a) Hard indicator --- pure-IS particle weighting.}
Replace the running cost $\ell$ by the indicator
$\ell^{\rm hard}(X_t, U_t) = \mathbb{1}[X_t \in \mathcal{B}_{\rm bad}(U_t)]$
and aggregate as a per-bin violation rate. Designed for pure-SMC$^2$
controllers that re-weight particles by their empirical violation
rate at each tempering step. Indicator-based; not differentiable in
$\thetactrl$.

\paragraph{(b) Soft sigmoid surrogate --- HMC compatible.}
Replace the indicator by a sigmoid surrogate
\begin{equation}
\ell^{\rm soft}(X_t, U_t) \;=\;
   \sigmoid\!\left(\,\frac{\beta_{\rm chance}\,\bigl(d_{\rm bad}(X_t, U_t)\bigr)}
                          {\rm scale}\,\right),
\label{eq:soft-surrogate}
\end{equation}
where $d_{\rm bad}(X_t, U_t)$ is a signed distance (positive when
$X_t$ is in the bad set, negative otherwise). The scaling parameter
$\beta_{\rm chance}$ controls the sharpness; as
$\beta_{\rm chance} \to \infty$ the surrogate converges to the hard
indicator. Differentiable in $\thetactrl$ via the trajectory; designed
for HMC inner kernels.

\paragraph{Empirical comparison: soft works, hard does not.}
Both variants have been implemented and tested side-by-side on the
chance-constrained control of a representative model. The soft
variant, run with HMC inside the outer SMC$^2$ and a $\beta$ annealed
from $\sim 50$ at low $\lambda$ to $\sim 10^4$ near $\lambda = 1$,
converges to the correct chance-constrained optimum and produces
schedules that respect the constraint at the prescribed rate. The
hard variant --- which matches the framework's importance-weighted
SMC$^2$ structure more naturally, but provides no gradient information
to drive HMC --- does \emph{not} converge to the same optimum in
practice. The empirical recommendation is therefore:

\begin{remark}[Use the soft variant]
For chance-constrained control problems, use the soft sigmoid
surrogate \eqref{eq:soft-surrogate} with HMC. The hard indicator
formulation is preserved in
\path{models/fsa_high_res/control_v5.py:evaluate_chance_constrained_cost_hard}
(in the FSA model dev sandbox repository) for documentation /
regression purposes only --- it is not the production cost function,
and should not be used to drive a deployed controller.
\end{remark}

The detailed numerical comparison and the per-bin behaviour
of both variants is in the FSA model documentation; that is where the
empirical run lives. The framework-level point for this document is
that both code paths exist, they target the same chance constraint
\eqref{eq:chance-constraint}, but only the soft path delivers a
working controller in practice.

\subsection{Implementation pointer}

\begin{itemize}
  \item \path{smc2fc/control/control_spec.py} --- the
        \texttt{ControlSpec} dataclass: cost callable, schedule
        decoder, horizon length, initial state, reference parameters.
  \item \path{smc2fc/control/tempered_smc_loop.py} --- the controller's
        outer loop. Built on the same SMC$^2$ engine as the inference
        side; the only difference is that the log-density it operates
        on is the cost-as-likelihood
        $-\beta\,J(\thetactrl)$ rather than $\log \Lhat_N(\thetadyn)$.
\end{itemize}
